\section{Related Work}
\label{sec:related_work}

\subsection{Text-to-SQL Benchmarks and Their Limitations}
The transition from academic-scale datasets (Spider, WikiSQL) to industry-representative benchmarks has been catalysed by the \highlight{BIRD} benchmark \parencite{li2024can}, which introduces 12,751 question--SQL pairs across 95 large-scale databases containing real-world dirty data. BIRD requires models to reason over complex multi-table joins, date format variations, and domain-specific value conventions. However, recent audits have revealed that up to 10\% of BIRD annotations contain semantic ambiguities or outright errors \parencite{floratou2024broken}, underscoring the need for systems that remain resilient to imperfect evaluation signals.

\subsection{Schema Pruning and Context Construction}
Providing an LLM with a full database schema---potentially hundreds of tables---exceeds context windows and degrades reasoning quality. \highlight{CHESS} \parencite{talaei2024chess} proposed a two-phase embedding pipeline:
\begin{enumerate}
    \item \textbf{Offline}: Index all table representations into a FAISS vector store \parencite{johnson2019faiss} using sentence-transformer embeddings \parencite{xiao2024bge}.
    \item \textbf{Online}: Embed only the user question and retrieve the Top-$k$ most similar tables via inner-product search, achieving sub-millisecond latency.
\end{enumerate}
Building on pruned schemas, \highlight{MCI-SQL} \parencite{zheng2024mcisql} demonstrated that enriching prompts with \emph{multiple facets} of context information---column data types, cardinality ratios (1:1 PK vs.\ 1:N FK), MIN/MAX statistics, and sample values---significantly improves schema linking accuracy. AgentSQL integrates both strategies sequentially: CHESS prunes the schema from $|\mathcal{S}|$ to $k$ tables, then MCI-SQL enriches the pruned subset with column-level metadata.

\subsection{Multi-Agent and Self-Correction Paradigms}
Multi-agent frameworks decompose the Text-to-SQL task into specialised sub-tasks. \highlight{MAC-SQL} \parencite{wang2024macsql} employs a Decomposer, a Selector, and a Refiner operating collaboratively. \highlight{MAGIC} \parencite{gu2024magic} generates self-correction guidelines from in-context exemplars, providing the LLM with an explicit error-checklist during generation. Execution-feedback approaches---\highlight{ReViSQL} \parencite{lee2024revisql} and \highlight{AGENTIQL} \parencite{xie2024agentiql}---iterate between generation and database execution to fix queries empirically.

AgentSQL distinguishes itself from these approaches along two axes:
\begin{itemize}
    \item \textbf{Asymmetric model allocation}: Unlike MAC-SQL's homogeneous agent pool, AgentSQL assigns a cost-efficient model (\texttt{llama-4-scout-17b}) for generation and reflection, reserving a high-reasoning model (\texttt{gemini-2.5-flash}) exclusively for the correction phase.
    \item \textbf{Error-type isolation}: Rather than sending all failures to a single corrector, a lightweight Reflector catches \emph{logical} mismatches (via back-translation) before execution, while a three-tier \texttt{SemanticErrorChecker} classifies \emph{syntactical} and \emph{semantic} failures post-execution---each with tailored correction suggestions.
\end{itemize}
